% This file should be rename with the speaker's last name. (e.g : smith)

% Write the author names in the order you wish them to appear
% and underline the one who will give the talk.

% Only the speaker's email address is required.

\begin{abstract_online}{Using CAS in the classroom: personal thoughts}{%
        \underline{Michel Beaudin}$^{1}$}{[michel.beaudin@etsmtl.ca]
        }{%
        $^1$ Service des enseignements généraux, École de technologie supérieure, Montréal (QC), Canada \newline}
    
My first contact with computer algebra goes back to \textbf{\emph{Derive}} in 1991. I will never forget what was written in the user manual introduction [1]: «Making mathematics more exciting and enjoyable is the driving force behind the development of the \textbf{\emph{Derive}} program.  The system is designed to eliminate the drudgery of performing long tedious mathematical calculations.  This gives you the freedom to explore different approaches to problems  \emph{--} approaches that you probably would not even consider if you had to do the calculations by hand ».  And this also applies, at different levels, to many other computer programs (CAS but also DGS for instance).
 
Textbooks such as [2] rarely give you the “freedom to explore different approaches to problems” but a daily use of computer algebra in the classroom may help to do so.  Since 1999, Texas Instruments CAS handheld and software technology are used in my engineering school by undergraduate students. Since then, I have decided to add in my teaching some aspects not covered or not enough exploited by textbooks.  So the talk will be about how this technology can be easily used to teach subjects where only pencil and paper techniques would discourage the user.  Chosen items in he following list will be used in the talk.
\begin{itemize}
 \item In algebra, we will take a look at a specific third degree polynomial equation where Nspire fails to do a good job.  With some help, satisfactory results will be obtained.
 \item In differential equations, we will see the importance of collecting similar terms when one is looking for a particular solution, especially if your CAS « deSolve » command is restricted to first and second order ODEs.
 \item Integral tables in calculus textbooks should be updated in order to benefit from computer algebra, nameley a better choice of antiderivatives and more explanations on the constant of integration. 
 \item Complex analysis  is a subject where computer algebra does not come in mind first but having access to a Laurent series function can change your opinion.
 \item Graphing facilities should be used anytime it can increase the understanding of concepts by students.  This is probably evident in calculus but this also applies to linear algebra!
\end{itemize}

Some among us thought computer algebra systems were going to change the way we teach mathematics.  It did but not as much as we should have expected. Blame for this whatever you want: the textbooks, the teachers, the students, the curriculum, the system.  But things are not so bad:  using CAS, many mathematics teachers have given themselves additional years of making "teaching more exciting and enjoyable".  



\textbf{Keywords} \newline{} Computer algebra systems, textbooks

\textbf{References} \newline{}
% Model for references to books.
[1] \emph{\textbf{DERIVE} User Manual, version 2}. Fourth edition, Soft Warehouse, March 1991. \newline{}
[2] \textsc{Stewart, J.}, \emph{Analyse, concepts et contextes, Fonctions de plusieurs variables}. 3ième édition, De Boeck, 2011. \newline{}
\end{abstract_online}
.
    